\section{Polynomial relations among Chern classes}
\label{section-relations-chern-classes}

\noindent
Let $(S, \delta)$ be as in Situation \ref{situation-setup}. Let $X$ be locally
of finite type over $S$. Let $\mathcal{E}_i$ be a finite collection of finite
locally free sheaves on $X$. By Lemma \ref{lemma-cap-commutative-chern}
we see that the Chern classes
$$
c_j(\mathcal{E}_i) \in A^*(X)
$$
generate a commutative (and even central) $\mathbf{Z}$-subalgebra of the
Chow cohomology algebra $A^*(X)$.
Thus we can say what it means for a polynomial in these Chern classes
to be zero, or for two polynomials to be the same. As an example, saying that
$c_1(\mathcal{E}_1)^5 + c_2(\mathcal{E}_2)c_3(\mathcal{E}_3) = 0$
means that the operations
$$
\CH_k(Y) \longrightarrow \CH_{k - 5}(Y), \quad
\alpha \longmapsto
c_1(\mathcal{E}_1)^5 \cap \alpha +
c_2(\mathcal{E}_2) \cap c_3(\mathcal{E}_3) \cap \alpha
$$
are zero for all morphisms $f : Y \to X$ which are locally of finite type.
By Lemma \ref{lemma-bivariant-zero}
this is equivalent to the requirement that given any morphism
$f : Y \to X$  where $Y$ is an integral scheme
locally of finite type over $S$ the cycle
$$
c_1(\mathcal{E}_1)^5 \cap [Y] +
c_2(\mathcal{E}_2) \cap c_3(\mathcal{E}_3) \cap [Y]
$$
is zero in $\CH_{\dim(Y) - 5}(Y)$.

\medskip\noindent
A specific example is the relation
$$
c_1(\mathcal{L} \otimes_{\mathcal{O}_X} \mathcal{N})
=
c_1(\mathcal{L}) + c_1(\mathcal{N})
$$
proved in Lemma \ref{lemma-c1-cap-additive}.
More generally, here is what happens when we tensor an
arbitrary locally free sheaf by an invertible sheaf.

\begin{lemma}
\label{lemma-chern-classes-E-tensor-L}
Let $(S, \delta)$ be as in Situation \ref{situation-setup}.
Let $X$ be locally of finite type over $S$.
Let $\mathcal{E}$ be a finite locally free sheaf of
rank $r$ on $X$. Let $\mathcal{L}$ be an invertible
sheaf on $X$. Then we have
\begin{equation}
\label{equation-twist}
c_i({\mathcal E} \otimes {\mathcal L})
=
\sum\nolimits_{j = 0}^i
\binom{r - i + j}{j} c_{i - j}({\mathcal E}) c_1({\mathcal L})^j
\end{equation}
in $A^*(X)$.
\end{lemma}

\begin{proof}
This should hold for any triple $(X, \mathcal{E}, \mathcal{L})$.
In particular it should hold when $X$ is integral and by
Lemma \ref{lemma-bivariant-zero}
it is enough to prove
it holds when capping with $[X]$ for such $X$. Thus assume
that $X$ is integral. Let $(\pi : P \to X, \mathcal{O}_P(1))$,
resp.\ $(\pi' : P' \to X, \mathcal{O}_{P'}(1))$ be the
projective space bundle associated to $\mathcal{E}$,
resp.\ $\mathcal{E} \otimes \mathcal{L}$. Consider the canonical morphism
$$
\xymatrix{
P \ar[rd]_\pi \ar[rr]_g & & P' \ar[ld]^{\pi'} \\
& X &
}
$$
see Constructions, Lemma \ref{constructions-lemma-twisting-and-proj}.
It has the property that
$g^*\mathcal{O}_{P'}(1)
= \mathcal{O}_P(1) \otimes \pi^* {\mathcal L}$.
This means that we have
$$
\sum\nolimits_{i = 0}^r
(-1)^i
(\xi + x)^i \cap \pi^*(c_{r - i}(\mathcal{E} \otimes \mathcal{L}) \cap [X])
=
0
$$
in $\CH_*(P)$, where $\xi$ represents
$c_1(\mathcal{O}_P(1))$ and $x$
represents $c_1(\pi^*\mathcal{L})$. By simple algebra this
is equivalent to
$$
\sum\nolimits_{i = 0}^r
(-1)^i \xi^i \left(
\sum\nolimits_{j = i}^r
(-1)^{j - i}
\binom{j}{i}
x^{j - i} \cap
\pi^*(c_{r - j}(\mathcal{E} \otimes \mathcal{L}) \cap [X])
\right)
=
0
$$
Comparing with
Equation (\ref{equation-chern-classes}) it follows from this that
$$
c_{r - i}(\mathcal{E}) \cap [X] =
\sum\nolimits_{j = i}^r
\binom{j}{i}
(-c_1(\mathcal{L}))^{j - i} \cap
c_{r - j}(\mathcal{E} \otimes \mathcal{L}) \cap [X]
$$
Reworking this (getting rid of minus signs, and renumbering) we get
the desired relation.
\end{proof}

\noindent
Some example cases of (\ref{equation-twist}) are
\begin{align*}
c_1(\mathcal{E} \otimes \mathcal{L})
& =
c_1(\mathcal{E}) +
r c_1(\mathcal{L}) \\
c_2(\mathcal{E} \otimes \mathcal{L})
& =
c_2(\mathcal{E}) +
(r - 1) c_1(\mathcal{E}) c_1(\mathcal{L}) +
\binom{r}{2} c_1(\mathcal{L})^2 \\
c_3(\mathcal{E} \otimes \mathcal{L})
& =
c_3(\mathcal{E}) +
(r - 2) c_2(\mathcal{E})c_1(\mathcal{L}) +
\binom{r - 1}{2} c_1(\mathcal{E})c_1(\mathcal{L})^2 +
\binom{r}{3} c_1(\mathcal{L})^3
\end{align*}








